\section{Keep complete vector groups separate from scalar tails}
The whole encoder calls its vector kernel while at least twelve input bytes remain,
advancing source by twelve and destination by sixteen. It then uses the scalar group
loop for the remaining zero through eleven bytes. Those bytes can include several
complete triples before the final padded group; a scalar tail is not necessarily just
one or two bytes.

The decoder excludes a final padded quartet from its vector-eligible count. It can then
call the sixteen-character kernel only for complete unpadded groups, advancing source
by sixteen and destination by twelve. Any remainder is handled by the validated scalar
quartet loop. A sixteen-character input ending in padding therefore uses the scalar
path in this implementation, since its unpadded prefix contains only twelve characters.

Keep original input count, remaining input, vector-eligible input and fixed total output
count distinct. The loop's source increment is not its output increment. Their units
are bytes in different representations, so using one stride for both corrupts the next
group's position. For exact in-place decoding, the output pointer falls progressively
behind unread input, which preserves the next group's source bytes.

Exercise lengths on both sides of the twelve-byte encoding boundary and the
sixteen-character decoding boundary. Include every input remainder modulo three,
along with padded text whose unpadded prefix is just shorter than a vector group.
A test of only complete large groups can leave an incorrect tail path untouched.
